% SPDX-License-Identifier: CC-BY-SA-4.0
% Author: Matthieu Perrin
% Part: <Nom de la partie>
% Section: <Nom de la section>
% Sub-section: <Nom de la sous-section>  % (facultatif, laisser vide si non utilisé)
% Frame: <Titre de la slide>

\begingroup

\begin{frame}{Rendu de scènes 3D}

  \on[top=-4mm, width=1.1\textwidth]{
    \Probleme{ray tracing}{
      \begin{itemize}
      \item Une description de scène géométrique 2D ou 3D
        \begin{itemize}
        \item Par exemple un ensemble de polygones plans
        \end{itemize}
      \item Une source de lumière ponctuelle \structure{\faLightbulbO}
      \item Une cible ponctuelle à éclairer \structure{\faCamera}
      \end{itemize}
    }{
      Existe-t-il une trajectoire de la source vers la cible ?
      \begin{itemize}
      \item La lumière est reflétée sur les objets de la scène
      \item On autorise les miroirs parfaits
      \end{itemize}
    }
  }
  
  \onBlock[y=-6mm, anchor=north]{Théorème (admis)}{
    \textsc{ray tracing} est indécidable.
  }


  \on[y=-24mm, x=25mm]{
    \begin{tikzpicture}[size=4mm]

      \draw[-latex, alert] (4.5,0.5) -- (0,5) -- (3,8) -- (4,7) -- (0,3) -- (2,1) -- (7,6) -- (8,5) ;
      \node[alert] at (8.5,4.5) {?};

      \draw[example, thick] (3.8,1) -- (0,1) -- (0,8) -- (4,8) -- (4,6) -- (5.8,6);
      \draw[example, thick] (6.2,6) -- (9,6) -- (9,1) -- (4.2,1);
      
      \node[example] at (5,0) {\faLightbulbO};
      \node[example] at (8,7.1) {\faCamera};
      \node[example] at (4.5,0.5) {$\bullet$};
      \node[example] at (7,7) {$\bullet$};
      
    \end{tikzpicture}
  }

  \footnoteref{Reif, Tygar, Yoshida. \emph{Computability and Complexity of Ray Tracing}. Discrete \& Computational Geometry (1994)}
  
\end{frame}

\endgroup
